% =================================================
% Справочный лист
% =================================================

\begin{center}
  {\fontsize{25}{30}\selectfont\bfseries\sffamily\color{Primary}
    Точки через равные промежутки}

  \vspace{0.3em}

  {\large\sffamily\color{GrayText}
    Как считать промежутки и находить расстояния}
\end{center}

\vspace{0.45em}

\section{Точки и промежутки --- не одно и то же}

\equalpointsbasic

\begin{propertybox}[Главное наблюдение]
  Между двумя соседними точками находится один промежуток. Если подряд
  отмечено несколько точек, то промежутков между первой и последней точками
  на один меньше, чем точек.
\end{propertybox}

\pointsvsintervalsfigure

Например, от первой точки до пятой нужно перейти четыре раза:
\[
  5-1=4.
\]

\begin{warningbox}[Типичная ошибка]
  Между первой и пятой точками не пять, а четыре промежутка. Номера точек
  показывают положения, а расстояние складывается из переходов между ними.
\end{warningbox}

\section{Промежутки между точками с заданными номерами}

\numberdifferencefigure

\begin{conclusionbox}[Правило]
  Чтобы найти число промежутков между точками с номерами, нужно из большего
  номера вычесть меньший:
  \[
    \text{число промежутков}
      =\text{больший номер}-\text{меньший номер}.
  \]
\end{conclusionbox}

Например, между 8-й и 27-й точками:
\[
  27-8=19
\]
промежутков.

\begin{notebox}[Проверка на небольших числах]
  Между 8-й и 9-й точками один промежуток: \(9-8=1\). Между 8-й и
  10-й --- два: \(10-8=2\). То же рассуждение работает и для больших
  номеров.
\end{notebox}

\section{Как найти расстояние}

Если соседние точки расположены на одинаковом расстоянии, то все промежутки
имеют одну и ту же длину.

\distanceexamplefigure

\begin{propertybox}[Расстояние между равноудалёнными точками]
  \[
    \text{расстояние}
      =\text{число промежутков}\cdot
       \text{длина одного промежутка}.
  \]
\end{propertybox}

В примере расстояние между 8-й и 27-й точками равно
\[
  (27-8)\cdot 7=19\cdot 7=133\text{ мм}.
\]

\section{Обратные задачи}

Иногда известно общее расстояние, но неизвестна длина одного промежутка.

\reverseexamplefigure

Между 3-й и 7-й точками четыре промежутка. Поэтому длина одного промежутка:
\[
  24:4=6\text{ см}.
\]

\begin{conclusionbox}[Если неизвестна длина одного промежутка]
  \[
    \text{длина одного промежутка}
      =\text{общее расстояние}:\text{число промежутков}.
  \]
\end{conclusionbox}

Можно встретить и задачу на номер точки.

\Needspace{15\baselineskip}
\begin{casebox}{Пример}
  Точки отмечены через \(4\) см. От 12-й точки вправо прошли \(20\) см.
  Какой номер имеет полученная точка?

  За \(20\) см пройдено
  \[
    20:4=5
  \]
  промежутков. Поэтому номер точки равен
  \[
    12+5=17.
  \]
\end{casebox}

Если двигаться в сторону уменьшения номеров, число промежутков нужно
вычитать из номера начальной точки.

\section{Столбы, деревья и другие объекты}

\postrowfigure

Задачи про точки часто маскируются под задачи про столбы, деревья, фонари,
остановки или отметки на линейке.

\begin{propertybox}[Объекты стоят в обоих концах ряда]
  Если в ряд поставлено \(n\) объектов, включая первый и последний, то между
  ними \(n-1\) одинаковых промежутков.
\end{propertybox}

Например, семь столбов образуют шесть промежутков. Если расстояние между
соседними столбами равно \(3\) м, то расстояние от первого до седьмого:
\[
  (7-1)\cdot3=18\text{ м}.
\]

\Needspace{10\baselineskip}
\begin{warningbox}[Сначала уточните, что считается]
  Вопрос «сколько столбов нужно на участке» и вопрос «каково расстояние от
  первого столба до последнего» устроены по-разному. Всегда нарисуйте хотя бы
  несколько объектов и отметьте промежутки между ними.
\end{warningbox}

\section{Чередующиеся объекты}

\alternatingtreesfigure

Если ряд начинается с сосны и деревья чередуются, то:
\begin{itemize}
  \item на нечётных местах стоят сосны;
  \item на чётных местах стоят берёзы.
\end{itemize}

Например, 21-е дерево --- сосна, а 34-е --- берёза.

\begin{questionsbox}[Проверьте себя]
  \begin{enumerate}
    \item Сколько промежутков между 4-й и 11-й точками?
    \item Почему между шестью подряд отмеченными точками пять промежутков?
    \item Как найти расстояние между 9-й и 25-й точками, если шаг равен
      \(6\) мм?
    \item Как найти длину одного промежутка по общему расстоянию?
    \item Сколько промежутков образуют 12 столбов, если столбы стоят в обоих
      концах ряда?
    \item Как определить вид дерева по номеру в чередующемся ряду?
  \end{enumerate}
\end{questionsbox}
